\begin{figure}[tbp]
\centering
\ViewHeading{(a) The full-torus seed leaves no room for strict growth}
\begin{tabular}{@{}llll@{}}
\toprule
Parameter & Relative quotient $H$ & Seed $V$ & Every admissible set of marks\\\midrule
$x\neq0$ & $\langle[v]\rangle$, dimension $1$ & $H$ & $\Nex=H$, no growth\\[3pt]
$x=0$ & $\C^d$, dimension $d$ & $0$ & $\Nex=0$, no growth\\\bottomrule
\end{tabular}
\medskip
\ViewHeading{(b) Different types of maps}
\begin{tikzpicture}
\node[vbox,text width=5.6cm] (scales) at (-3.55,0)
 {\textbf{Arithmetic scale map}\\
  $T_{q,k}:\mathcal C_x(kq)_1\to\mathcal C_x(q)_1$\\
  Different scales and powered coefficients};
\node[vresult,text width=5.6cm] (marks) at (3.55,0)
 {\textbf{Response endomorphism}\\$\overline G_j:H\to H$\\One supplied fixed quotient};
\node[vinput,text width=11.9cm] (extra) at (0,-2)
 {Missing additional construction: quotient identifications or a new finite model,\\
  with proved boundary preservation and intertwining};
\draw[vconditional] (scales.south)--(extra.north west);
\draw[vconditional] (extra.north east)--(marks.south);
\node[font=\footnotesize,align=center] at (0,-3.35)
 {Dashed arrows are requirements, not maps constructed by the papers.\\
  Trivial powered characters can also change the relative dimension.};
\end{tikzpicture}
\caption{The arithmetic torus does not automatically provide a growing
response module. Its full degree-one cocycle equations force the top
table; substituting the wedge changes the datum. The lower comparison
keeps cross-scale chain maps separate from fixed-quotient marks.}
\label{fig:response-torus}
\end{figure}